\section{Exact certificate and claim ledger}
\label{sec:tpc45-certificate}

The reproducibility directory contains
\begin{center}
 \path{experiments/tpc45_certificate.py}.
\end{center}
It uses the Python standard library, exact Gaussian-integer arithmetic,
exhaustive finite sign cubes, and \texttt{fractions.Fraction}.  A
successful run writes
\path{experiments/tpc45_certificate.json}.  The checked families are:
\begin{enumerate}[label=\textup{(\roman*)},leftmargin=2.4em]
 \item the actual toy labels \(u=d(\ell d j+1)\), the packet-defined
 band, uniqueness of a band prime, exact physical sign deletion, and
 equality of the original and grouped energies on every remaining sign
 environment;
 \item the grouped diagonal inequality and equality in the abstract
 fiber-sum operator norm;
 \item the physical-low affine form on every high-prime environment;
 \item the low-averaged pure Gram projection and its all-minus synthesis
 identity on exhaustive low and high cubes;
 \item the fixed-determinant ladder, its affine parameter step, and the
 common physical sign on an equal-residual fiber; and
 \item every rational exponent and endpoint margin used in the paper.
\end{enumerate}
The current certificate performs \(32894\) exact checks.  Its canonical
payload hash is
\begin{center}
\texttt{e9c3bdf309863cf560fd31e5253e1429236281069101df26c7fe689d6faa438e}.
\end{center}
The hash covers the report before the hash field itself is appended.

The certificate is deliberately finite.  It checks identities and
arithmetic ledgers, not the assertions that \(\tau(n)=X^{o(1)}\), that
primes exist in each fixed power subinterval, or that a density-one
probability tends to one.  Those are proved analytically in the text.
Most importantly, it does not sample or certify M\"obius cancellation,
a fixed-determinant Bessel estimate, or a prime-pair asymptotic.

For clarity, the logical status of the main statements is as follows.
\begin{enumerate}[label=\textup{\arabic*.},leftmargin=2.2em]
 \item The actual coefficient expansion, prime-incidence alternatives,
 cofactor deletion, fiber bounds, grouped diagonal bounds, affine Gram
 forms, and determinant-ladder parameterization are unconditional finite
 or asymptotic theorems under the inherited TPC packet hypotheses.
 \item The mesoscopic physical-band freezing and its two-parameter
 band/cutoff tradeoff are unconditional consequences of those theorems
 and the robust-face argument.
 \item The larger cutoff \((\log X)^{1003/1000}\) applies to the
 shortened touched component.  It is not claimed for the untouched part
 or the full packet.
 \item The fixed-determinant synthesis estimate in
 \cref{sec:tpc45-next-gate} is a proposed next input, not a theorem of
 this paper.
 \item The comparisons in \cref{sec:tpc45-literature} are applicability
 audits.  A mismatch with a cited theorem is not an impossibility result
 for every future method.
 \item No statement here implies the twin-prime conjecture, a
 Hardy--Littlewood asymptotic, a breach of the sieve parity barrier, or
 the complete deterministic TPC endpoint.
\end{enumerate}

Thus the positive result is a real partial de-randomization of actual
TPC coefficients, while the remaining gap is displayed rather than
silently promoted to an assumption.
